\section*{अध्याय \ref{ch:g5:flower-to-fruit} --- फूल से फल तक}

\begin{solution}{exo:g5:flower-to-fruit:1}
\omterm{def:g3:flowers-fruits-seeds:parts}{पुंकेसर} \omterm{def:g4:animal-reproduction:sexual}{नर} हैं: उनके पराग-कण \omterm{def:g4:animal-reproduction:sexual}{नर} जननिक \omterm{def:g5:first-look-at-cells:cell}{कोशिकाएँ} लिए रहते हैं।
\omterm{def:g3:flowers-fruits-seeds:parts}{स्त्रीकेसर} \omterm{def:g4:animal-reproduction:sexual}{मादा} है: उसके भीतर \omterm{prop:g5:flower-to-fruit:roles}{बीजांड} पड़े रहते हैं, और हर \omterm{prop:g5:flower-to-fruit:roles}{बीजांड} में
एक \omterm{def:g4:animal-reproduction:sexual}{अंडाणु} होता है।
\end{solution}

\begin{solution}{exo:g5:flower-to-fruit:2}
\omterm{def:g5:flower-to-fruit:steps}{परागण}: \omterm{def:g3:flowers-fruits-seeds:parts}{पराग} का एक कण किसी \omterm{def:g3:flowers-fruits-seeds:parts}{पुंकेसर} से किसी \omterm{def:g3:flowers-fruits-seeds:parts}{स्त्रीकेसर} की चोटी तक सफ़र
करता है। \omterm{def:g5:flower-to-fruit:steps}{निषेचन}: उस कण की नली \omterm{def:g4:animal-reproduction:sexual}{नर} \omterm{def:g5:first-look-at-cells:cell}{कोशिका} को नीचे किसी \omterm{prop:g5:flower-to-fruit:roles}{बीजांड} तक ले
जाती है, जहाँ वह \omterm{def:g4:animal-reproduction:sexual}{अंडाणु} से जुड़ जाती है। \omterm{def:g5:flower-to-fruit:steps}{परागण} पहले आता है ---
मुलाक़ात से पहले यात्रा।
\end{solution}

\begin{solution}{exo:g5:flower-to-fruit:3}
किसी \omterm{def:g3:flowers-fruits-seeds:parts}{पुंकेसर} में बना; \omterm{prop:g3:sorting-living-things:groups}{कीट} या हवा के साथ किसी \omterm{def:g3:flowers-fruits-seeds:parts}{स्त्रीकेसर} की चोटी तक
पहुँचाया गया --- यानी \omterm{def:g5:flower-to-fruit:steps}{परागण}; \omterm{def:g3:flowers-fruits-seeds:parts}{स्त्रीकेसर} के भीतर नीचे तक एक महीन नली
उगाई; और उसकी \omterm{def:g4:animal-reproduction:sexual}{नर} \omterm{def:g5:first-look-at-cells:cell}{कोशिका} किसी \omterm{prop:g5:flower-to-fruit:roles}{बीजांड} तक पहुँचकर \omterm{def:g4:animal-reproduction:sexual}{अंडाणु} से जुड़ी ---
यानी \omterm{def:g5:flower-to-fruit:steps}{निषेचन}।
\end{solution}

\begin{solution}{exo:g5:flower-to-fruit:4}
हर निषेचित \omterm{prop:g5:flower-to-fruit:roles}{बीजांड} एक \omterm{def:g2:from-seed-to-plant:seed}{बीज} बन जाता है; \omterm{def:g3:flowers-fruits-seeds:parts}{स्त्रीकेसर} फूलकर \omterm{prop:g3:flowers-fruits-seeds:transform}{फल} बन जाता है;
और \omterm{def:g3:flowers-fruits-seeds:parts}{पंखुड़ियाँ} (तथा \omterm{def:g3:flowers-fruits-seeds:parts}{पुंकेसर}) मुरझाकर गिर जाते हैं।
\end{solution}

\begin{solution}{exo:g5:flower-to-fruit:5}
हर दाना एक निषेचित \omterm{prop:g5:flower-to-fruit:roles}{बीजांड} है जो \omterm{def:g2:from-seed-to-plant:seed}{बीज} बन गया। ख़ाली कक्षों का मतलब है बिना
\omterm{def:g5:flower-to-fruit:steps}{निषेचन} वाले \omterm{prop:g5:flower-to-fruit:roles}{बीजांड} --- बहुत कम \omterm{def:g3:flowers-fruits-seeds:parts}{पराग} उतरा था, और \omterm{prop:g3:flowers-fruits-seeds:transform}{फल} उस कमी के चारों ओर
असमान फूला।
\end{solution}

\begin{solution}{exo:g5:flower-to-fruit:6}
उनका \omterm{def:g3:flowers-fruits-seeds:parts}{पराग} हवा पर सवार होता है: धूल जितना महीन, बहुत बड़ी मात्रा में बना,
और हवा में उछाल दिया गया। इसका दाम है बरबादी --- लगभग हर कण कहीं
बेकार जगह उतरता है, इसलिए सटीकता की जगह मात्रा को लेनी पड़ती है।
\end{solution}

\begin{solution}{exo:g5:flower-to-fruit:7}
मुलाक़ात नहीं, तो \omterm{def:g3:flowers-fruits-seeds:parts}{स्त्रीकेसर} पर \omterm{def:g3:flowers-fruits-seeds:parts}{पराग} नहीं --- यानी पहला क़दम, \omterm{def:g5:flower-to-fruit:steps}{परागण},
हुआ ही नहीं, इसलिए न कोई नली उगी और न कोई \omterm{prop:g5:flower-to-fruit:roles}{बीजांड} निषेचित हुआ; और \omterm{def:g3:flowers-fruits-seeds:parts}{फूल}
पूरा मुरझा जाता है।
\end{solution}

\begin{solution}{exo:g5:flower-to-fruit:8}
उत्तर खुला है। उम्मीद यह है: हाथ से परागित, धागे से चिह्नित फूलों में
साफ़ तौर पर ज़्यादा बार \omterm{prop:g3:flowers-fruits-seeds:transform}{फल} लगता है --- कूची ने मधुमक्खी का काम किया।
\end{solution}

\begin{solution}{exo:g5:flower-to-fruit:9}
\omterm{def:g4:animal-reproduction:sexual}{नर} \omterm{def:g5:first-look-at-cells:cell}{कोशिकाएँ}: \omterm{def:g3:flowers-fruits-seeds:parts}{पुंकेसरों} के पराग-कणों में। \omterm{def:g4:animal-reproduction:sexual}{अंडाणु}: \omterm{def:g3:flowers-fruits-seeds:parts}{स्त्रीकेसर} के भीतर के
\omterm{prop:g5:flower-to-fruit:roles}{बीजांडों} में। \omterm{def:g5:flower-to-fruit:steps}{निषेचन}: पराग-नली के तल पर, \omterm{prop:g5:flower-to-fruit:roles}{बीजांड} के भीतर। नए व्यक्ति की
पहली \omterm{def:g5:first-look-at-cells:cell}{कोशिका}: \omterm{prop:g5:flower-to-fruit:roles}{बीजांड} में पड़ा निषेचित \omterm{def:g4:animal-reproduction:sexual}{अंडाणु} --- जो बढ़कर \omterm{def:g2:from-seed-to-plant:seed}{बीज} का नन्हा
\omterm{def:g1:plants-around-us:plant}{पौधा} बनता है, ठीक वैसे ही जैसे कोई \omterm{def:g1:animals-around-us:animal}{जंतु} अपनी पहली \omterm{def:g5:first-look-at-cells:cell}{कोशिका} से बढ़ता
है।
\end{solution}

\begin{solution}{exo:g5:flower-to-fruit:10}
फ़सल ख़राब रहेगी --- ज़्यादातर \omterm{def:g3:flowers-fruits-seeds:parts}{फूल} बिना \omterm{prop:g3:flowers-fruits-seeds:transform}{फल} के ही गिर गए। मशीनरी दुरुस्त
थी और फूलों को कोई चोट नहीं आई; जो नाकाम हुआ वह पहला क़दम था, यानी
\omterm{def:g5:flower-to-fruit:steps}{परागण}: दो निर्णायक हफ़्तों तक हरकारे घर बैठे रहे, इसलिए ज़्यादातर
\omterm{def:g3:flowers-fruits-seeds:parts}{स्त्रीकेसरों} को \omterm{def:g3:flowers-fruits-seeds:parts}{पराग} मिला ही नहीं और ज़्यादातर \omterm{prop:g5:flower-to-fruit:roles}{बीजांड} कभी निषेचित नहीं
हुए।
\end{solution}

\begin{solution}{exo:g5:flower-to-fruit:11}
भँवरे \omterm{def:g5:flower-to-fruit:steps}{परागण} बेचते हैं --- यानी ढुलाई का क़दम: ऐसी इमारत के भीतर
\omterm{def:g3:flowers-fruits-seeds:parts}{पुंकेसरों} से \omterm{def:g3:flowers-fruits-seeds:parts}{स्त्रीकेसरों} तक \omterm{def:g3:flowers-fruits-seeds:parts}{पराग} पहुँचाना जहाँ बाहर का कोई परागणकर्ता
घुस ही नहीं सकता। टमाटर का \omterm{def:g3:flowers-fruits-seeds:parts}{पराग} हिलाने भर से आसानी से झड़ जाता है,
इसलिए पंखे का कंपन और चलती हवा हर \omterm{def:g3:flowers-fruits-seeds:parts}{फूल} के भीतर मोटे तौर पर वह ढुलाई कर
देते हैं --- यानी एक यांत्रिक हरकारा। चेरी के \omterm{def:g3:flowers-fruits-seeds:parts}{पराग} को \omterm{def:g1:plants-around-us:tree}{पेड़} से \omterm{def:g1:plants-around-us:tree}{पेड़} तक,
\omterm{def:g3:flowers-fruits-seeds:parts}{फूल} से \omterm{def:g3:flowers-fruits-seeds:parts}{फूल} तक जाना पड़ता है, और उतनी सटीकता सिर्फ़ वही शरीर देता है जो
बारी-बारी फूलों पर बैठता हो; हॉल की हवा निशाना नहीं साध सकती, इसलिए
चेरी के लिए जीवित हरकारे का कोई सस्ता विकल्प नहीं है।
\end{solution}
